% GENERATED FILE. Edit canonical lessons, then run npm run generate:books.
% canonical-source-hash: fnv1a64:6646c9105d55b401
% canonical-lessons: UR-C241-palna, UR-C241-bona, UR-C241-ugana, UR-C241-jharu-lagana, UR-C241-ponchhna

\chapter{To bring up, To sow, To cultivate, To sweep, To wipe}
\label{ch:ur-to-bring-up-to-sow-to-cultivate-to-sweep-to-wipe}
\hlchaptermodality{241}

\begin{chapteropening}
\textbf{By the end of this chapter:} \emph{I can say to bring up, to sow, to cultivate, to sweep and to wipe.}

Complete the last lesson of chapter 241: I can say to bring up, to sow, to cultivate, to sweep and to wipe.
\end{chapteropening}

\section[\texorpdfstring{\ur{پالنا} (pālnā)}{پالنا (pālnā)}]{\ur{پالنا} (pālnā) — to bring up, to keep (animals)}

\label{lesson:UR-C241-palna}



\begin{quote}
Before the new one: say the Urdu for to halt, then the Urdu for to turn.
\end{quote}

\subsection*{You'll want to know: \ur{پالنا}}
\textbf{\ur{پالنا}} — \emph{pālnā} — \textquotedblleft{}to bring up, to keep (animals)\textquotedblright{}.

\begin{cousinweb}
\textbf{\ur{پالنا}} (\emph{pālnā}) is a tadbhava, a word inherited from Sanskrit and reshaped by sound change. It comes from \emph{pālayati}, \textquotedblleft{}guards, protects\textquotedblright{} (CDIAL 8129).
\end{cousinweb}

\begin{grammarlens}[title={What you've built}]
One more word for this chapter.
\end{grammarlens}

\subsection*{Your turn}
\begin{itemize}
\raggedright
  \item \emph{Say it:} \emph{pālnā}
  \item \emph{Say it:} \emph{pālnā}, once more
  \item \emph{Say it:} say \emph{pālnā}
  \item \emph{Recall:} say the Urdu for to halt, then the Urdu for to turn, then say \emph{pālnā} again
\end{itemize}

\begin{tcolorbox}[breakable,colback=teal!4,colframe=teal!35!black,title={Before you move on}]
What does \ur{پالنا} mean? (\textquotedblleft{}To bring up, to keep (animals)\textquotedblright{}.) Say it once more.
\end{tcolorbox}
\section[\texorpdfstring{\ur{بونا} (bonā)}{بونا (bonā)}]{\ur{بونا} (bonā) — to sow}

\label{lesson:UR-C241-bona}



\begin{quote}
Before the new one: say the Urdu for to turn, then the Urdu for to bring up.
\end{quote}

\subsection*{You'll want to know: \ur{بونا}}
\textbf{\ur{بونا}} — \emph{bonā} — \textquotedblleft{}to sow\textquotedblright{}.

\begin{cousinweb}
\textbf{\ur{بونا}} (\emph{bonā}) goes back to Sanskrit \emph{vapati}, \textquotedblleft{}scatters, sows\textquotedblright{} (CDIAL 11282). CDIAL marks the form \emph{bonā} as taken over from Gujarati.
\end{cousinweb}

\begin{grammarlens}[title={What you've built}]
One more word for this chapter.
\end{grammarlens}

\subsection*{Your turn}
\begin{itemize}
\raggedright
  \item \emph{Say it:} \emph{bonā}
  \item \emph{Say it:} \emph{bonā}, once more
  \item \emph{Say it:} say \emph{bonā}
  \item \emph{Recall:} say the Urdu for to turn, then the Urdu for to bring up, then say \emph{bonā} again
\end{itemize}

\begin{tcolorbox}[breakable,colback=teal!4,colframe=teal!35!black,title={Before you move on}]
What does \ur{بونا} mean? (\textquotedblleft{}To sow\textquotedblright{}.) Say it once more.
\end{tcolorbox}
\section[\texorpdfstring{\ur{اگانا} (ugānā)}{اگانا (ugānā)}]{\ur{اگانا} (ugānā) — to cultivate, to grow (plants)}

\label{lesson:UR-C241-ugana}



\begin{quote}
Before the new one: say the Urdu for to bring up, then the Urdu for to sow.
\end{quote}

\subsection*{You'll want to know: \ur{اگانا}}
\textbf{\ur{اگانا}} — \emph{ugānā} — \textquotedblleft{}to cultivate, to grow (plants)\textquotedblright{}.

\begin{cousinweb}
\textbf{\ur{اگانا}} (\emph{ugānā}) is a tadbhava, a word inherited and reshaped by sound change. It comes from \emph{*udgāti}, \textquotedblleft{}rises, comes forth\textquotedblright{}, a form CDIAL reconstructs from its descendants (CDIAL 1954). \emph{ugānā} is its causing form, \textquotedblleft{}to make spring up\textquotedblright{}.
\end{cousinweb}

\begin{grammarlens}[title={What you've built}]
One more word for this chapter.
\end{grammarlens}

\subsection*{Your turn}
\begin{itemize}
\raggedright
  \item \emph{Say it:} \emph{ugānā}
  \item \emph{Say it:} \emph{ugānā}, once more
  \item \emph{Say it:} say \emph{ugānā}
  \item \emph{Recall:} say the Urdu for to bring up, then the Urdu for to sow, then say \emph{ugānā} again
\end{itemize}

\begin{tcolorbox}[breakable,colback=teal!4,colframe=teal!35!black,title={Before you move on}]
What does \ur{اگانا} mean? (\textquotedblleft{}To cultivate, to grow (plants)\textquotedblright{}.) Say it once more.
\end{tcolorbox}
\section[\texorpdfstring{\ur{جھاڑو} \ur{لگانا} (jhāṛū lagānā)}{جھاڑو لگانا (jhāṛū lagānā)}]{\ur{جھاڑو} \ur{لگانا} (jhāṛū lagānā) — to sweep}

\label{lesson:UR-C241-jharu-lagana}



\begin{quote}
Before the new one: say the Urdu for to sow, then the Urdu for to cultivate.
\end{quote}

\subsection*{You'll want to know: \ur{جھاڑو} \ur{لگانا}}
\textbf{\ur{جھاڑو} \ur{لگانا}} — \emph{jhāṛū lagānā} — \textquotedblleft{}to sweep\textquotedblright{}.

\begin{cousinweb}
\textbf{\ur{جھاڑو} \ur{لگانا}} (\emph{jhāṛū lagānā}) is \textquotedblleft{}to apply the broom\textquotedblright{}. \emph{jhāṛū} is built on \emph{jhāṛnā}, \textquotedblleft{}to sweep\textquotedblright{}, which CDIAL lists under \emph{*jhāṭayati}, \textquotedblleft{}shakes down\textquotedblright{} (CDIAL 5328); \emph{lagānā} is the causing form of \emph{lagnā}, \textquotedblleft{}to be attached\textquotedblright{}, which CDIAL lists under Sanskrit \emph{lagyati} (CDIAL 10895).
\end{cousinweb}

\begin{grammarlens}[title={What you've built}]
One more word for this chapter.
\end{grammarlens}

\subsection*{Your turn}
\begin{itemize}
\raggedright
  \item \emph{Say it:} \emph{jhāṛū lagānā}
  \item \emph{Say it:} \emph{jhāṛū lagānā}, once more
  \item \emph{Say it:} say \emph{jhāṛū lagānā}
  \item \emph{Recall:} say the Urdu for to sow, then the Urdu for to cultivate, then say \emph{jhāṛū lagānā} again
\end{itemize}

\begin{tcolorbox}[breakable,colback=teal!4,colframe=teal!35!black,title={Before you move on}]
What does \ur{جھاڑو} \ur{لگانا} mean? (\textquotedblleft{}To sweep\textquotedblright{}.) Say it once more.
\end{tcolorbox}
\section[\texorpdfstring{\ur{پونچھنا} (ponchhnā)}{پونچھنا (ponchhnā)}]{\ur{پونچھنا} (ponchhnā) — to wipe}

\label{lesson:UR-C241-ponchhna}



\begin{quote}
Before the new one: say the Urdu for to cultivate, then the Urdu for to sweep.
\end{quote}

\subsection*{You'll want to know: \ur{پونچھنا}}
\textbf{\ur{پونچھنا}} — \emph{ponchhnā} — \textquotedblleft{}to wipe\textquotedblright{}.

\begin{cousinweb}
\textbf{\ur{پونچھنا}} (\emph{ponchhnā}) is a tadbhava, a word inherited from Sanskrit and reshaped by sound change. It comes from \emph{proñchati}, \textquotedblleft{}wipes off\textquotedblright{} (CDIAL 9011).
\end{cousinweb}

\begin{grammarlens}[title={What you've built}]
One more word for this chapter.
\end{grammarlens}

\subsection*{Your turn}
\begin{itemize}
\raggedright
  \item \emph{Say it:} \emph{ponchhnā}
  \item \emph{Say it:} \emph{ponchhnā}, once more
  \item \emph{Say it:} say \emph{ponchhnā}
  \item \emph{Recall:} say the Urdu for to cultivate, then the Urdu for to sweep, then say \emph{ponchhnā} again
\end{itemize}

\begin{tcolorbox}[breakable,colback=teal!4,colframe=teal!35!black,title={Before you move on}]
What does \ur{پونچھنا} mean? (\textquotedblleft{}To wipe\textquotedblright{}.) Say it once more.
\end{tcolorbox}
